% title: LEETCODE 564 (Solution)
% overview: This document describes a method to solve a hard solution on leetcode (564. The method is to find the next palindrome and the previous palindrome, then compare the distance between the two palindromes and the input number.)
% date: 2024-09
% tags: Math, Palindrome, Leetcode, C++, Problem solving

\documentclass{article}
\usepackage[utf8]{inputenc}
\usepackage[T1]{fontenc}
\usepackage[french]{babel}
\usepackage{amsmath}
\usepackage{graphicx} % Required for inserting images
\usepackage{minted}   % Compiler avec -shell-escape (activé par défaut sur Overleaf)

\title{Fanazavana}
\author{Rakotoarisoa Ando}
\date{July 2024}

\begin{document}

\maketitle

\section{564. Find the Closest Palindrome [Hard - Daily]}

\textbf{Input} : $n$ une chaîne représentant un nombre entre $0$ et $10^{18}$, soit $N = |n|$.

\textbf{Problème} : Trouver un palindrome $P$ différent de $n$ qui minimise la distance entre $P$ et $n$ (en terme de valeur absolue ilay resahina eto).

\textbf{Solution} :

\textbf{Observation 1} : Il faut chercher le palindrome akaiky indrindra, c'est-à-dire le prochain palindrome et le précédent palindrome.

\textbf{Observation 2} : Si on peut juste transformer la chaîne en palindrome et mahafeno an'ireo fepetra ireo ilay izy dia on est bon. Il suffit d'assigner $s_{N - i + 1} \leftarrow s_i$.

\textbf{Observation 3} : Sinon il faut se salir un peu les mains. Sans perte de généralité, considérons l'algorithme qui va trouver le prochain palindrome $P$ différent de $n$. On supposera alors que $P > n$. On le veut le plus proche possible de $n$, donc quel indice (puissance de $10$) changer ?

Si on choisit un indice $i < j < \frac{N}{2}$, alors en choisissant de modifier $i$ plutôt que $j$, on aura
\[
P \leftarrow P + a\,10^{i - 1} + b\,10^{N - i}
\]
tandis que si on choisit $j$, on aura
\[
P \leftarrow P + c\,10^{j - 1} + d\,10^{N - j}
\]
tel que $a\,10^{i - 1} + b\,10^{N - i} > c\,10^{j - 1} + d\,10^{N - j}$. Donc si on choisit $i$, on « s'éloigne » plus de $n$ que si l'on choisit $j$.

On peut donc procéder à un algorithme glouton en parcourant d'abord le plus au centre (meilleur local), puis en s'étendant vers les bords. Il est à noter qu'il faut aussi gérer le cas particulier $9\ldots9$ ($N$ fois), pour lequel la réponse est $1\underbrace{0\ldots0}_{N-1}1$.

\begin{minted}{c++}
using ll = long long;

class Solution {
public:
    string nextPalindrome(string s, string& ref)
    {
        int n = s.size();
        for (int i = 0; i < n / 2; i++)
            if (s[i] != s[n - i - 1])
                s[n - i - 1] = s[i];
        if (stoll(s) > stoll(ref))
            return (s);
        int i = n / 2 - (n + 1) % 2;
        while (i >= 0 && s[i] == '9' && s[n - i - 1] == '9')
        {
            s[i] = '0', s[n - i - 1] = '0';
            i--;
        }
        if (i < 0) s += '1', s[++i] = '1';
        else s[n - i - 1] = ++s[i];
        return (s);
    }

    string previousPalindrome(string s, string& ref)
    {
        int n = s.size();
        for (int i = 0; i < n / 2; i++)
            if (s[i] != s[n - i - 1])
                s[n - i - 1] = s[i];
        if (stoll(s) < stoll(ref))
            return (s);
        int i = n / 2 - (n + 1) % 2;
        while (i >= 0 && s[i] == '0' && s[n - i - 1] == '0')
        {
            s[i] = '9', s[n - i - 1] = '9';
            i--;
        }
        if (i < 0) s.pop_back();
        else
        {
            if (!i && s[i] == '1' && n > 1) s = string(n - 1, '9');
            else s[n - i - 1] = --s[i];
        }
        return (s);
    }

    ll distance(string s, string t)
    {
        return abs(stoll(s) - stoll(t));
    }

    string nearestPalindromic(string n) {
        string nxt = nextPalindrome(n, n);
        string prv = previousPalindrome(n, n);
        ll d1 = distance(nxt, n), d2 = distance(prv, n);
        if (d1 == d2) return prv;
        return (d1 < d2) ? nxt : prv;
    }
};
\end{minted}

\end{document}